\section{The isotrivial case when the second Betti number is four}
\label{iso:section}
\label{sec:isotrivial}

\begin{proposition}\label{iso:exclusion}
Let \(f:X\to B\) be the projective Lagrangian fibration of
Section~\ref{sec:setup}.  If \(b_2(X)=4\) and the polarized period map
of its smooth fibers has zero variation, then every germ of \(B\) is a
quotient singularity.
\end{proposition}

Throughout the section we assume these hypotheses.  The proof starts
with the original fibration \(f\).  Its constant abelian fiber and
rational transcendental plane produce a normal variety \(T\) with
two finite maps: a Galois cover \(p:T\to B\), and a map \(a:T\to A\)
to an abelian variety.  We compare their ramification indices with
the multiplicity of a divisor in a fiber of \(f\).  The resulting
inequalities control both the linear monodromy and the lifting of
vector fields near a minimal stratum.  This is the isotrivial cover;
the finite cover of \(B\) by an abelian variety constructed from a
second Lagrangian fibration in the next section has a different role.

Finite-cover descriptions and the Hodge theory of
isotrivial Lagrangian systems are developed systematically by
Kim--Laza--Martin \cite[Section~2.4]{KimLazaMartin2026}; the construction
below retains the two distinct ramification maps needed for our
rank-four argument.

\begin{lemma}\label{iso:vertical-divisor}
The Picard number of \(X\) is two, and
\[
 \mathbb T_{\Q}:=\NS(X)_{\Q}^{\perp}
 \quad\text{satisfies}\quad
 \mathbb T_{\Q}\otimes_{\Q}\C
   =\C[\sigma]\oplus\C[\overline{\sigma}].
\]
For every prime divisor \(D\subset B\), there is a unique prime divisor
\(F_D\subset X\) dominating \(D\).  Near the generic point of \(D\),
\[
                         f^*D=m_DF_D
\]
for a positive integer \(m_D\).
\end{lemma}

\begin{proof}
Since \(h^{2,0}(X)=h^{0,2}(X)=1\), we have \(h^{1,1}(X)=2\).
An ample class \(h\) and the nonzero isotropic class \(e\) are linearly
independent, so they span \(\NS(X)_{\Q}\).  Their Beauville form is
nondegenerate, and its orthogonal complement is the stated rational
plane.

Let \(F\) be any divisor dominating \(D\).  Its restriction to a general
smooth fiber is zero, whereas \(h\) restricts nontrivially and \(e\)
restricts to zero.  Therefore \([F]=a e\) with \(a\in\Q\).
Intersecting with an ample class to the power \(2n-1\) shows \(a>0\).
Choose positive integers \(k,l\) with \(k[F]=l e\).
After multiplying once more if necessary, this is equality of integral
classes.  Since \(H^1(X,\cO_X)=0\), the exponential sequence gives
\[
                 \cO_X(kF)\simeq f^*H^{\otimes l}.
\]
The section defining \(kF\) is pulled back from a section of
\(H^{\otimes l}\), because \(f_*\cO_X=\cO_B\).
Its zero set consequently contains every component over the generic
point of \(D\).  Its support is \(F\), so there is only one such
component.  A local uniformizer of \(D\) now has divisor \(m_DF_D\).
\end{proof}

\begin{lemma}\label{iso:abelian-model}
There exist a normal projective variety \(T\), a finite group \(G\),
an \(n\)-dimensional abelian variety \(P\), and an \(n\)-dimensional
abelian variety \(A\), with the following properties.
\begin{enumerate}
\item The map \(p:T\to B\) is finite Galois with group \(G\), and the
main base change of \(X\) is birational over \(T\) to \(P\times T\).
The action on the first factor is by constant affine automorphisms
preserving an ample line bundle on \(P\).
\item There is a finite surjective \(G\)-equivariant morphism
\(a:T\to A\), where \(G\) acts affinely on \(A\).
The linear actions on \(P\) and \(A\) have the same kernel.
\item Writing this kernel as \(K\), and writing \(A_l\) for the quotient
of \(A\) by the finite group of translations induced by \(K\), there
are finite morphisms
\[
       T_l:=T/K \xrightarrow{\ a_l\ } A_l,
       \qquad
       T_l\xrightarrow{\ p_l\ }B.
\]
The group \(G_l:=G/K\) is the Galois group of \(p_l\), and acts on
\(A_l\) with faithful linear part.
\end{enumerate}
\end{lemma}

\begin{proof}
Choose an ample line bundle on \(X\), and use its restriction to polarize
the smooth fibers.  Constancy of the polarized moduli map means that
the geometric generic fiber, after choosing an origin, is a fixed
polarized abelian variety \(P\).
The origin and a polarized isomorphism are defined over a finite
extension of \(\C(B)\): each is specified by finitely many algebraic
data over its algebraic closure.  Alternatively, a sufficiently high
level structure gives this isomorphism on a dense open by the fine
moduli space of polarized abelian varieties.

We can also make the actual line bundle constant.  Once its numerical
polarization has been matched with a fixed ample bundle \(L_P\) on
\(P\), their difference lies in \(\Pic^0(P)\).
The polarization isogeny
\[
 P\longrightarrow\Pic^0(P),\qquad
 x\longmapsto t_x^*L_P\otimes L_P^{-1}
\]
allows us to remove that difference by a translation after a further
finite extension.  Take a Galois closure and normalize \(B\) in it.
The resulting \(T\) is projective and normal.
The descent automorphisms of \(P\) preserve \(L_P\).
The group \(\Aut(P,L_P)\) is finite: its translation subgroup is the
finite kernel of the polarization isogeny, and its origin-preserving
subgroup is a discrete subgroup of the compact unitary group of the
polarization.  Thus each descent automorphism, initially over the
function field of \(T\), is a constant affine automorphism over \(\C\).
This proves the first assertion.

Take a \(G\)-equivariant projective resolution
\(\widetilde T\to T\).
The rational map \(P\times\widetilde T\dashrightarrow X\) pulls
\(\sigma\) back to a holomorphic two-form \(\widetilde\sigma\) on
\(P\times\widetilde T\).  To see regularity, resolve the graph; its
pulled-back form descends across the birational map to the smooth
product, since holomorphic differential forms are unchanged by a
proper birational modification of smooth varieties.
Choose a basis \(\alpha_1,\ldots,\alpha_n\) of \(H^0(P,\Omega_P^1)\).
The K\"unneth decomposition and the Lagrangian condition give
\begin{equation}\label{iso:cross-term}
 \widetilde\sigma
    =\sum_{i=1}^n\alpha_i\wedge\gamma_i+\beta,
 \qquad
 \gamma_i\in H^0(\widetilde T,\Omega_{\widetilde T}^1),\quad
 \beta\in H^0(\widetilde T,\Omega_{\widetilde T}^2).
\end{equation}
There is no two-form term from \(P\).
At a general point the pulled-back form is nondegenerate and \(P\)
is Lagrangian; hence the \(\gamma_i\) are pointwise independent there.

We spell out the rationality needed for an abelian quotient.
Graph pullback and projection onto the K\"unneth summand give a
morphism of rational Hodge structures
\[
 \mathbb T_{\Q}\longrightarrow
 H^1(P,\Q)\otimes H^1(\widetilde T,\Q).
\]
Contract its image with \(H^1(P,\Q)^\vee\), and denote the resulting
rational subspace of \(H^1(\widetilde T,\Q)\) by \(W_{\Q}\).
By Lemma~\ref{iso:vertical-divisor} and
\eqref{iso:cross-term},
\[
 W_{\Q}\otimes\C
   =\langle\gamma_1,\ldots,\gamma_n\rangle_{\C}
       \oplus
     \langle\overline{\gamma_1},\ldots,
                 \overline{\gamma_n}\rangle_{\C}.
\]
It is a rational weight-one Hodge substructure of dimension \(2n\).
To form the corresponding Albanese quotient, take the annihilator
of \(W_{\Q}\) in
\(H_1(\operatorname{Alb}(\widetilde T),\Q)\).
It is a rational Hodge subspace, so its real span and its lattice
define a connected complex subtorus.  Restriction of an ample
polarization makes this an abelian subvariety.
The quotient is an abelian variety \(A\) of dimension \(n\);
its one-forms pull back to the span of the \(\gamma_i\).
The induced morphism
\(\widetilde a:\widetilde T\to A\) is generically finite.
The subspace \(W_{\Q}\) is \(G\)-stable; the Albanese functor consequently
gives an affine action on \(A\) making \(\widetilde a\) equivariant.

We next descend this morphism to \(T\).
Let \(C'\) be a curve contracted by \(\widetilde T\to T\).
The rational map from \(P\times C'\) to \(X\), interpreted on a resolved
graph, has image in one fiber of \(f\).
The pullback of \(\sigma\) on that graph is zero.  Here the Lagrangian
condition also applies if the image lies in a smaller subvariety of a
fiber component: resolve the component, on which the pulled-back form
vanishes, and dominate the smaller subvariety by proper base change.
Restriction of \eqref{iso:cross-term} and contraction in the \(P\)
directions now show that each \(\gamma_i\) vanishes on the normalization
of \(C'\).  It follows that \(\widetilde a\) is constant on every
contracted curve.  It is constant on each connected resolution fiber
as well: a positive-dimensional projective image of such a fiber in
an abelian variety would contain the image of a noncontracted curve.
The fibers of \(\widetilde T\to T\) are connected and \(T\) is normal,
so this proves descent to a morphism \(a:T\to A\).

Equivariance gives a generically finite morphism
\(\bar a:B=T/G\to A/G\).
Pull back an ample line bundle from the projective quotient \(A/G\).
The result is a nef and big line bundle on \(B\); because
\(\rho(B)=1\), it is numerically a positive multiple of an ample
class and hence is ample.  Thus \(\bar a\) contracts no curve and is
finite.  The commutative quotient square then shows that \(a\) is
finite too.

Finally the cross pairing in \eqref{iso:cross-term} is perfect between
the tangent spaces of the two abelian factors.  Its invariance shows
that an element has trivial linear action on \(P\) if and only if it
has trivial linear action on \(A\).  Such an element acts on \(A\)
by a translation.  Quotienting its finite translation image gives
\(A_l\), and equivariance gives the stated finite map from \(T/K\).
The derivative representation is unchanged by an isogeny, so the
linear action of \(G_l\) on \(A_l\) is faithful.
\end{proof}

\begin{lemma}\label{iso:ramification}
Let \(D\subset B\) be a prime divisor, let \(E\subset T\) be a prime
divisor over \(D\), and put \(D_A=a(E)\).  Write
\[
                 M=e_E(p),\qquad L=e_E(a)
\]
for the ramification indices of the two finite maps, and let \(m_D\)
be the fiber multiplicity in Lemma~\ref{iso:vertical-divisor}.
Then
\begin{equation}\label{iso:multiplicity-bound}
                             Lm_D\le M.
\end{equation}
The linear part of a generator of the cyclic inertia is either the
identity or a pseudoreflection, on both \(A\) and \(P\).
If \(M_l,L_l\) are the analogous indices for \(p_l,a_l\), then
\begin{equation}\label{iso:linear-ramification-bound}
                              L_l\le M_l.
\end{equation}
\end{lemma}

\begin{proof}
The maps and divisors being compared fit into the following diagram.
The upper square commutes as rational maps, and the lower square is
the quotient square from Lemma~\ref{iso:abelian-model}.
\begin{center}
\begin{tikzpicture}[>=Latex, every node/.style={font=\small}]
 \node (prod) at (0,1.6) {\(P\times T\supset P\times E\)};
 \node (source) at (6.1,1.6) {\(X\supset F_D\)};
 \node (cover) at (0,0) {\(T\supset E\)};
 \node (base) at (6.1,0) {\(B\supset D\)};
 \node (abelian) at (0,-1.6) {\(A\supset D_A\)};
 \node (quotient) at (6.1,-1.6) {\(A/G\)};
 \draw[->,dashed] (prod) -- node[above] {\(\phi\)} (source);
 \draw[->] (prod) -- node[left] {\(\operatorname{pr}_T\)} (cover);
 \draw[->] (source) -- node[right,align=left]
   {\(f\)\\\(\operatorname{ord}_{F_D}(f^*t)=m_D\)} (base);
 \draw[->] (cover) -- node[above,align=center]
   {\(p\) finite Galois\\\(e_E(p)=M\)} (base);
 \draw[->] (cover) -- node[left,align=right]
   {\(a\) finite\\\(e_E(a)=L\)} (abelian);
 \draw[->] (base) -- node[right] {\(\bar a\)} (quotient);
 \draw[->] (abelian) -- node[above] {quotient} (quotient);
\end{tikzpicture}
\end{center}
Here \(t\) is a local uniformizer of \(D\), and \(\phi\) is induced by
the birational identification over \(T\).  The integers \(M\) and
\(L\) are local indices along \(E\), not the degrees of the finite
maps.  They compare to the multiplicity \(m_D\) on the original
source \(X\), which may be greater than one.

All assertions concern generic points of divisors.  We may therefore
discard smaller subsets, split the unramified tangential extension,
and work with smooth transverse polydiscs.  Choose a transverse
parameter \(z\) for \(E\), absorbing a unit so that \(p^*t=z^M\).
Cyclic inertia fixes \(E\) pointwise.  Since \(a\) is finite and
equivariant, it fixes the divisor \(D_A\) pointwise.
A finite affine automorphism with a fixed point is conjugate by a
translation to its linear part; fixing a divisor makes its linear
part either the identity or a pseudoreflection.
The perfect cross pairing gives the same assertion on \(P\).

Let \(\omega=\sigma^n/n!\), with an irrelevant nonzero scalar absorbed
into the choice of a translation-invariant volume form on \(P\).
In the top exterior power of \eqref{iso:cross-term}, every term must
contain all \(n\) \(P\)-directions.  Thus the term \(\beta\) makes no
contribution, and on the smooth product model
\[
             \phi^*\omega=\operatorname{vol}_P\wedge a^*
                                      \operatorname{vol}_A.
\]
Its order along \(P\times E\) is \(L-1\).  More precisely, near the
generic point of \(E\), it is \(u z^{L-1}\operatorname{vol}_{\mathrm{prod}}\),
where
\(\operatorname{vol}_{\mathrm{prod}}
=\operatorname{vol}_P\wedge\operatorname{vol}_T\)
for a nowhere vanishing local top form \(\operatorname{vol}_T\) on \(T\),
and \(u\) is pulled back from a unit at the generic point of \(E\).

Let \(v=\operatorname{ord}_{F_D}\) be the normalized valuation of
the divisor in Lemma~\ref{iso:vertical-divisor}.  The product field is
the compositum \(\C(P\times T)=\C(X)\C(T)\).  Geometric integrality
of the generic fiber makes this a Galois extension of \(\C(X)\) with
group \(G\).  Extend \(v\) to a
normalized divisorial valuation \(w\) of the product field
\(\C(P\times T)\), choosing the extension above \(E\).  Indeed,
conjugating any extension by \(G\) preserves its restriction \(v\)
and moves its base center through all divisors of \(T\) over \(D\).
Write \(r\) for the ramification index of \(w/v\).  Since
\(v(t)=m_D\),
\begin{equation}\label{iso:valuation-multiplicity}
                         r m_D=Mw(z).
\end{equation}
The center of \(w\) on the proper product projects dominantly to \(E\).
In particular the product is smooth at the generic point of this
center, and the unit \(u\) above is a \(w\)-unit.
Denote its log discrepancy by \(A_{\mathrm{prod}}(w)\).
The volume form on the smooth \(X\) has order zero at \(v\).
The order of the same pulled-back top form can therefore be computed
in the finite extension of valuation rings and on the product model:
\begin{equation}\label{iso:top-form-valuation}
 \begin{aligned}
 \operatorname{ord}_w(\phi^*\omega)
   &=r\,\operatorname{ord}_v(\omega)+(r-1)=r-1\\
   &=(L-1)w(z)+A_{\mathrm{prod}}(w)-1.
 \end{aligned}
\end{equation}
The first line uses the different exponent \(r-1\) of a tame
extension of discrete valuation rings; the residue extension is
separable in characteristic zero.  The second line uses
\(\operatorname{ord}_w(\operatorname{vol}_{\mathrm{prod}})
=A_{\mathrm{prod}}(w)-1\).
On the smooth product, the pair defined by the smooth divisor
\(z=0\) is log canonical, so \(A_{\mathrm{prod}}(w)\ge w(z)\).
Consequently \eqref{iso:top-form-valuation} gives
\begin{equation}\label{iso:valuation-discrepancy}
 r=(L-1)w(z)+A_{\mathrm{prod}}(w)
                       \ \ge\ Lw(z).
\end{equation}
This last inequality can also be read in local coordinates:
the pullback of a form containing \(dz\) has order at least
\(w(z)-1\).
Combining \eqref{iso:valuation-multiplicity} and
\eqref{iso:valuation-discrepancy} proves
\eqref{iso:multiplicity-bound}.

Let \(k\) be the inertia order of \(T\to T_l\) along this divisor.
Its image on \(A\) is a group of translations fixing \(a(z=0)\),
so that image is trivial.  Quotienting \(A\) by the translation
group is an isogeny and is everywhere unramified.  The indices in
the two compositions consequently give
\[
                          M=kM_l,\qquad L=kL_l.
\]
In particular \(L_lm_D\le M_l\), which implies
\eqref{iso:linear-ramification-bound}.
\end{proof}

\begin{lemma}\label{iso:order-two}
Let \(g_D\) be a local inertia generator at a prime base divisor \(D\).
The linear part of \(g_D^{m_D}\) on \(P\) does not have order two.
\end{lemma}

\begin{proof}
Set \(m=m_D\), and make the base change \(t=s^m\).
Near the generic point of \(D\), the normalized source \(X'\) of this
base change is smooth and its map to \(X\) is \'{e}tale.
Indeed \(f^*t\) has the single divisor \(mF_D\).
In the regular local rings of \(X\) one can write
\(f^*t=u h^m\), with \(u\) a unit; adjoining \(s\) and normalizing
adjoins an \(m\)-th root of the unit.  This is \'{e}tale.
Equivalently, the extension is unramified in codimension one and
purity applies.  The divisor \(s=0\) on \(X'\) is reduced.
These statements hold on a neighborhood of the whole fiber over
the generic point of \(D\), which is the neighborhood used below.

Put \(\delta=\gcd(M,m)\).
On a component of the normalized compositum of the two base covers,
use a parameter \(z'\) so that
\[
 z=(z')^{m/\delta},\qquad
 s=(z')^{M/\delta}.
\]
Its cyclic inertia over the rooted base has order
\[
             M'=\frac M\delta,
 \qquad\text{and its generator acts on \(P\) as }g_D^m.
\]
The order of the pulled-back volume on the new product is \(L'-1\),
where
\begin{equation}\label{iso:rooted-indices}
                         L'=\frac{Lm}{\delta}.
\end{equation}
This follows by substituting \(z=(z')^{m/\delta}\) in
\(z^{L-1}dz\).  In particular
\(0<L'\le M'\), by \eqref{iso:multiplicity-bound}.

Write \(T'\) for the chosen normalized cover, \(E'\) for its divisor
\(z'=0\), and \(\phi':P\times T'\dashrightarrow X'\) for the induced
rational map.  Let \(\omega'\) be the pullback of \(\omega\) to \(X'\);
it is nowhere vanishing on the smooth neighborhood just described.
We first show that the full affine inertia action has a common fixed
point.  Take an actual divisor of \(X'\) over \(s=0\), with normalized
valuation \(v'\), and extend it to a normalized divisorial valuation
\(w'\) of the product field above \(E'\).
Geometric integrality of the generic fiber makes this field extension
Galois with the same group as \(T'\) over the rooted base.
Let \(r'=e(w'/v')\); in this tame Galois extension it is also the
inertia order.  Since \(v'(s)=1\), reducedness gives
\[
                     r'=M'w'(z').
\]
The center \(Z\) of \(w'\) on the proper product projects dominantly
to \(E'\).  Thus
\[
             \kappa(E')\ \subset\ \kappa(Z)\ \subset\ \kappa(w').
\]
Every inertia element fixes \(\kappa(w')\); its restriction to the
base therefore fixes \(\kappa(E')\) and belongs to the cyclic base
inertia of order \(M'\).
Hence \(r'\le M'\), and therefore
\[
                         r'=M',\qquad w'(z')=1.
\]
Thus the whole cyclic base inertia fixes \(\kappa(w')\), and hence
\(\kappa(Z)\).  The generic point of \(Z\) supplies a point of \(P\)
over \(\kappa(Z)\) fixed by this full affine group.  Its common fixed
locus is a constant closed subscheme of \(P\), so nonemptiness after
that field extension implies nonemptiness over \(\C\).
Translation by a complex fixed point conjugates the affine action to
the linear action.  In particular their orders agree.

Suppose now that the linear action has order two.
It is a pseudoreflection, so its determinant on the top fiber form
is \(-1\).  A generator rotates \(z'\) by a primitive \(M'\)-th
root \(\zeta\).
The leading coefficient of the volume is independent of the
\(P\)-variable, and inertia fixes the tangential base residue field.
Invariance of the volume therefore gives
\[
                          -\zeta^{L'}=1.
\]
Thus \(M'\) is even and \(L'\equiv M'/2\pmod {M'}\).
The bound \(0<L'\le M'\) makes this
\[
                              L'=M'/2.
\]

Consider instead the product divisor \(P\times E'\) itself, with its
normalized valuation \(w_0\).
The cyclic base inertia acts on \(\kappa(w_0)\) through its affine
action on \(P\), of order two.  If \(v_0\) is the normalized restriction
to the field of \(X'\), every element of the inertia of \(w_0/v_0\)
fixes \(\kappa(E')\) and therefore lies in this base inertia.
It is exactly the kernel of the displayed residue action, so
\[
                         e_0=e(w_0/v_0)=M'/2.
\]
Properness over the rooted base supplies a center for \(v_0\) on
\(X'\).  It lies in the smooth neighborhood of the whole fiber over
the generic point of \(s=0\).  Thus, even before knowing whether
\(v_0\) is a divisor on \(X'\), the nowhere vanishing form
\(\omega'\) satisfies
\(\operatorname{ord}_{v_0}(\omega')=A_{X'}(v_0)-1\).
Computing the top-form order and the order of \(s\) in the finite
extension gives
\[
 \begin{aligned}
 L'-1
   &=e_0\operatorname{ord}_{v_0}(\omega')+(e_0-1)
     =e_0A_{X'}(v_0)-1,\\
 M'&=w_0(s)=e_0v_0(s).
 \end{aligned}
\]
Since \(L'=e_0=M'/2\), these equalities say
\[
                  A_{X'}(v_0)=1,\qquad v_0(s)=2.
\]
A normalized divisorial valuation over a smooth variety
with log discrepancy one is the valuation of a divisor on that
variety: a center of codimension at least two has log discrepancy
at least two, as is seen by taking regular parameters at its
generic point.  Hence \(v_0\) is an actual divisor of the smooth
source \(X'\).
Its displayed multiplicity two contradicts the reduced divisor
defined by \(s\).  This proves the lemma.
\end{proof}

The two local conclusions have separate uses.  The inequality
\(L_l\le M_l\) will make pulled-back differentials regular near the
minimal stratum.  The order bound for \(g_D^{m_D}\), together with
normal generation by powered meridians, now bounds the dimension.

\begin{lemma}\label{iso:dimension-bound}
The linear monodromy \(G_l\) acts irreducibly on
\(\operatorname{Lie}(P)\), is generated by pseudoreflections of order
greater than two, and satisfies \(n\le4\).
\end{lemma}

\begin{proof}
Choose a connected dense open \(B^\circ\) on which \(f\) is smooth,
the base cover is unramified, and the product description in
Lemma~\ref{iso:abelian-model} holds.  Write \(X^\circ=f^{-1}(B^\circ)\).
The smooth proper map \(X^\circ\to B^\circ\) has connected fibers,
so its map of fundamental groups is surjective.
Since \(X\) is simply connected,
\(\pi_1(X^\circ)\) is normally generated by small meridians around
the divisorial boundary of \(X^\circ\).  One obtains this statement
by filling loops and putting the filling discs in general position;
the subset of complex codimension at least two does not contribute.
Equidimensionality ensures that the inverse image of a subset of
base codimension at least two also has codimension at least two.
By Lemma~\ref{iso:vertical-divisor}, a meridian around \(F_D\)
projects to the \(m_D\)-th power of a meridian around \(D\).

The unramified restriction of \(T\to B\) is connected, and hence its
monodromy has full image \(G\).
Translation parts act trivially on \(H^1(P)\); the linear monodromy
is therefore exactly \(G_l\).
The preceding normal generation says that it is generated by
conjugates of the linear parts of \(g_D^{m_D}\).
Lemma~\ref{iso:ramification} makes each nontrivial such part a
pseudoreflection; Lemma~\ref{iso:order-two} excludes order two.
Notice that generation uses the powered meridians.  Replacing
them by unpowered inertia would lose this conclusion.

The global invariant cycle theorem, applied to the smooth projective
family \(X^\circ\to B^\circ\) and its smooth projective compactification
\(X\), identifies
\[
       H^2(P,\C)^{G_l}
       =\operatorname{im}\bigl(H^2(X,\C)\longrightarrow H^2(P,\C)\bigr).
\]
The hypotheses here are the algebraic smooth family over the connected
open and the proper smooth compactification; no assertion about
invariants on an arbitrary analytic disc is being used
\cite[Th\'eor\`eme~4.1.1(ii)]{Deligne1971}.
The classes \(e,[\sigma],[\overline\sigma]\) restrict to zero,
and the ample class restricts nontrivially.
These four classes span \(H^2(X,\C)\), so the displayed space has
dimension one.

If \(\operatorname{Lie}(P)\) had a nontrivial invariant splitting,
average a positive Hermitian form over the finite group and take an
orthogonal invariant splitting.  The imaginary parts of its
restrictions to the two summands give two independent invariant
real alternating forms on the underlying real tangent space.
Their constant forms represent independent classes in
\(H^2(P,\C)^{G_l}\), a contradiction.  Thus the representation is
irreducible.

We use the following precise consequence of the reflection
classification.  An irreducible finite complex reflection group
generated by reflections of order greater than two has dimension
at most four.  Indeed, in the imprimitive case the imprimitivity
spaces are lines and every reflection is either diagonal or has
a two-coordinate block
\[
                       \begin{pmatrix}0&u\\u^{-1}&0\end{pmatrix}.
\]
The latter has order two; diagonal reflections alone cannot generate
an irreducible representation of dimension greater than one.
These are exactly the alternatives in
\cite[\S2, pp.~276--277]{ShephardTodd1954}.
In the primitive case,
\cite[Proposition~(5.9)(iii), p.~430]{Cohen1976}
states that dimension at least five permits no reflection of
order greater than two.  Applying this consequence proves \(n\le4\).
\end{proof}

Suppose, towards a contradiction, that a nonquotient germ remains.
Take \(S_0,\cS,d,c\) from Proposition~\ref{str:minimum}.
The transverse klt germ on a maximal chart has dimension \(c\).
It cannot have dimension one, since a normal curve is smooth.
It cannot have dimension two, since a klt surface germ is a quotient
germ and its maximal quasi-\'etale cover is smooth
\cite[Proposition~(6.11)]{ClemensKollarMori1988}.
Thus \(c\ge3\).
Proposition~\ref{eu:euler-zero} excludes \(d=0\).
Lemma~\ref{iso:dimension-bound} now forces
\begin{equation}\label{iso:elliptic-case}
              n=4,\qquad c=3,\qquad d=1,\qquad
              \chi_{\mathrm{top}}(S_0)=0.
\end{equation}
The coarse curve \(S_0\) is normal and projective, hence smooth.
Its Euler number makes it a curve of genus one.
The remainder of the proof rules out this last possibility.

\begin{lemma}\label{iso:local-fields}
For \(z_0\in T_l\), put \(J=\operatorname{Stab}_{G_l}(z_0)\) and
\[
                V_J=\operatorname{Lie}(A_l)^J.
\]
Every translation field belonging to \(V_J\) lifts to a holomorphic
vector field on a neighborhood of \(z_0\) in \(T_l\).  This lift is
\(J\)-invariant and descends through \(p_l\).
If \(p_l(z_0)\in S_0\), then \(\dim V_J\le1\).
\end{lemma}

\begin{proof}
First describe the branch divisors of \(a_l\).
If \(R\subset T_l\) has ramification index \(L_l>1\), then
\eqref{iso:linear-ramification-bound} implies \(M_l>1\) for \(p_l\).
Its nontrivial cyclic inertia fixes \(R\) pointwise and therefore
fixes the divisor \(a_l(R)\) pointwise.
Since the linear action of \(G_l\) is faithful, a generator acts
on \(A_l\) by an affine pseudoreflection.
Consequently \(a_l(R)\) is a component of its fixed locus:
it is a translate of an abelian subvariety of codimension one.
We call such a translate an abelian hyperplane.

After shrinking about \(z_0\), every ramification divisor that meets
the neighborhood contains \(z_0\).
Its inertia is contained in \(J\), since an automorphism fixing
the divisor pointwise fixes \(z_0\).
The tangent hyperplane of its image therefore contains \(V_J\).

Let \(v\in V_J\).
The finite separable map \(a_l\) gives a unique meromorphic lift
\(\widetilde v\) of the constant field \(v\).
It is regular in codimension one.
At an unramified generic point this is immediate.
At a general point of a ramification divisor choose coordinates
in which \(a_l\) is
\[
        (z,y_2,\ldots,y_n)\longmapsto
                     (z^{L_l},y_2,\ldots,y_n).
\]
Since \(v\) is tangent to the abelian branch hyperplane, its normal
component is zero and the lift has no pole.
Normality then extends the derivation everywhere: for any regular
function \(h\), the meromorphic function \(\widetilde v(h)\)
belongs to every local height-one ring and thus to the normal
local ring itself.
Uniqueness of the lift and \(J\)-invariance of \(v\) show that
\(\widetilde v\) is \(J\)-invariant.
The germ of \(p_l\) is the quotient by \(J\), so the field descends
to that germ of \(B\).

These descended fields have independent evaluations.
To verify this without assuming either germ smooth, center
exponential coordinates of \(A_l\) at \(a_l(z_0)\), so the finite
affine \(J\)-action becomes linear.
Choose \(J\)-invariant linear coordinate functions
\(\ell_1,\ldots,\ell_r\), dual to a basis
\(v_1,\ldots,v_r\) of \(V_J\).
The invariant functions \(a_l^*\ell_i\) descend to functions
\(u_i\) on the base germ, and the descended fields \(\xi_j\) satisfy
\[
                              \xi_j(u_i)=\delta_{ij}.
\]
They therefore split \(r\) smooth directions by their local flows.
The same directions persist on the maximal quasi-\'etale chart:
the fields lift uniquely across a quasi-\'etale finite map
(check regularity in codimension one and use normality), and
the pulled-back functions retain these identities.
On \(S_0\) the evaluation rank on that chart is \(d=1\), by
Proposition~\ref{str:minimum}.  Thus \(r\le1\).
\end{proof}

\begin{lemma}\label{iso:regular-differentials}
For every local holomorphic function \(u\) on \(B\), its pulled-back
differential has an expression
\begin{equation}\label{iso:differential-coefficients}
                 p_l^*(du)=\sum_{i=1}^n h_i\,a_l^*\alpha_i,
\end{equation}
where the \(\alpha_i\) are a basis of translation-invariant
one-forms on \(A_l\) and the \(h_i\) are holomorphic on the
corresponding neighborhood of \(T_l\).
At a point \(z_0\) with stabilizer \(J\), the covector
\(\sum_i h_i(z_0)\alpha_i\) is \(J\)-invariant.
\end{lemma}

\begin{proof}
The coefficients are uniquely defined meromorphic functions, since
\(a_l\) is generically \'{e}tale.
It is enough to check regularity at every prime divisor of \(T_l\).
There is nothing to prove where \(a_l\) is unramified.
At a general ramification point write \(M=M_l\), \(L=L_l\).
Use linear coordinates \((x,y_2,\ldots,y_n)\) on \(A_l\) adapted
to the affine fixed hyperplane \(x=0\).
The inertia fixes the \(y_i\) and acts on \(x\) by a root of unity.
The restrictions of \(a_l^*y_i\) are generically independent along
the divisor.  They can be used as tangential coordinates on
\(T_l\), and, after absorbing a unit in the normal parameter,
\[
                       a_l^*x=z^L.
\]
Choose \(z\) so that the cyclic inertia rotates it by a primitive
\(M\)-th root.  The local invariant ring of this cyclic quotient
has coordinates \(z^M,y_2,\ldots,y_n\).
Thus \(p_l^*u=F(z^M,y)\) for a holomorphic function \(F\), and
\[
 d(p_l^*u)
   =\frac ML z^{M-L}F_t(z^M,y)\,d(z^L)
       +\sum_{i=2}^n F_{y_i}(z^M,y)\,dy_i.
\]
All coefficients are holomorphic because \(L\le M\).
The displayed coordinate covectors are constant linear combinations
of the \(\alpha_i\), so this proves the required codimension-one
regularity.  Normality extends every \(h_i\).

Expression \eqref{iso:differential-coefficients} is unique on the
function field.  Its left side is \(J\)-invariant and its right
side transforms by the linear representation on the \(\alpha_i\).
Evaluation at the \(J\)-fixed point \(z_0\) therefore gives an
invariant covector, as asserted.
\end{proof}

\begin{lemma}\label{iso:elliptic-lift}
Let \(C\) be an irreducible component of
\((p_l^{-1}(S_0))_{\mathrm{red}}\).
Then \(C\) is a smooth curve of genus one.
Its image \(E=a_l(C)\) is an elliptic translate in \(A_l\), and
\(a_l|_C:C\to E\) and \(p_l|_C:C\to S_0\) are \'{e}tale.
At every \(z\in C\),
\[
       \operatorname{Lie}(A_l)^{\operatorname{Stab}_{G_l}(z)}
                            =\operatorname{Lie}(E).
\]
Here \(\operatorname{Lie}(E)\) means the constant tangent direction
of the translate \(E\).
\end{lemma}

\begin{proof}
All components \(C\) are projective curves mapping finitely onto
\(S_0\).  Let \(J_0\) be their generic pointwise stabilizer.
More explicitly, discard the proper fixed loci of those finitely
many group elements that do not fix \(C\) identically.
On the remaining open the stabilizer is the subgroup \(J_0\)
fixing \(C\) pointwise.
Equivariance implies
\[
                         a_l(C)\subset A_l^{J_0}.
\]
Every component of this fixed locus is a translate of an abelian
subvariety with tangent space \(V_{J_0}\).
Lemma~\ref{iso:local-fields} bounds its dimension by one, whereas
finiteness of \(a_l\) makes the image of \(C\) one-dimensional.
It follows that the image is exactly an elliptic translate \(E\),
and \(V_{J_0}=\operatorname{Lie}(E)\).

Write \(\nu:\widehat C\to C\) for the normalization and let
\[
           b:\widehat C\to S_0,\qquad
           a_C:\widehat C\to E
\]
be the two finite morphisms.
At any \(q\in\widehat C\), choose a parameter \(u\) on the smooth
curve \(S_0\) near \(b(q)\), and extend it to a local function on
\(B\).  Such an extension exists because \(S_0\) has its reduced
closed-subspace structure.
Restricting \eqref{iso:differential-coefficients} gives
\begin{equation}\label{iso:curve-differential}
                          b^*du=\eta\,a_C^*\theta
\end{equation}
near \(q\), where \(\theta\) is a nonzero translation-invariant
one-form on \(E\) and \(\eta\) is holomorphic.
There is no assumption here that \(T_l\) is smooth along \(C\).
To justify restriction, first pull
\eqref{iso:differential-coefficients} to a resolution of \(T_l\).
Both sides are regular ordinary differentials there and agree
generically, hence everywhere.  A component of the proper base
change dominating \(\widehat C\) then gives the stated equality
on \(\widehat C\), since pullback of differentials along a dominant
map to a smooth curve in characteristic zero is injective.

Equation \eqref{iso:curve-differential} implies pointwise
\[
                       \operatorname{ram}_q(b)
                          \ge\operatorname{ram}_q(a_C).
\]
Both targets have genus one.  Riemann--Hurwitz therefore gives
the same total ramification degree, \(2g(\widehat C)-2\), for
both maps.  The pointwise inequalities must all be equalities.
In particular \(\eta(q)\ne0\).

Put \(z=\nu(q)\) and \(J=\operatorname{Stab}_{G_l}(z)\).
The covector in Lemma~\ref{iso:regular-differentials} is
\(J\)-invariant, and the nonvanishing of \(\eta(q)\) says that
its restriction to \(\operatorname{Lie}(E)\) is nonzero.
Since \(J\) contains \(J_0\), we have
\[
                         V_J\subset V_{J_0}
                                      =\operatorname{Lie}(E).
\]
A finite-dimensional representation of a finite group in
characteristic zero has equally dimensional invariant vector
and covector spaces, by the averaging projector.
Thus \(V_J\ne0\), and the displayed inclusion is equality.
This holds at every point and on every normalization branch.

Fix a nonzero \(v_0\in\operatorname{Lie}(E)\).
Its unique meromorphic lift through \(a_l\) is holomorphic near
each point of \(C\), by Lemma~\ref{iso:local-fields}, and the
lifts agree on overlaps by uniqueness.
Call the resulting field \(\widetilde v_0\).
It has nonzero evaluation everywhere on \(C\), since
\(da_l(\widetilde v_0)=v_0\).
Locally it descends through \(p_l\) to a field on \(B\).
Its flow preserves the intrinsically defined minimum stratum
\(S_0\), and hence its lifted flow preserves the reduced inverse
image and each of its irreducible curve components.
Thus \(\widetilde v_0\) is tangent to \(C\).

A reduced curve carrying such a nonvanishing ambient holomorphic
field is smooth: choose a function whose derivative under the
field is nonzero, and use the local flow to split off its disc
direction.  An invariant reduced curve in this product is the
disc times a reduced zero-dimensional germ.
It follows that \(C=\widehat C\) is smooth.
Since \(da_l(\widetilde v_0)=v_0\ne0\), the map \(C\to E\) is
unramified.  It is consequently an \'{e}tale cover of an elliptic
curve, so \(g(C)=1\).
The equality of ramification divisors already proved also makes
\(C\to S_0\) unramified.
\end{proof}

\begin{proof}[Proof of Proposition~\ref{iso:exclusion}]
If a nonquotient germ remained, \eqref{iso:elliptic-case} and
Lemma~\ref{iso:elliptic-lift} would supply \(C\) and a nonzero
translation direction \(v_0\) as above.
Choose a positive integer \(r\) such that
\(\mathcal L=\cO_B(rK_B)\) is a line bundle, and put
\[
                         \mathcal M=p_l^*\mathcal L.
\]
We construct a holomorphic connection on \(\mathcal M|_C\);
in particular we must check both its extension and its gluing.

Near \(z\in C\), the field \(\widetilde v_0\) is invariant under
the point stabilizer and descends to a holomorphic field
\(\xi_z\) on the corresponding base germ.
On \(B_{\mathrm{reg}}\), Lie differentiation by \(\xi_z\) acts
on \(r\)-fold canonical tensors.  It extends to an operator
\[
             \mathcal D_z:\mathcal L\longrightarrow\mathcal L.
\]
Indeed \(\mathcal L\) is the reflexive extension of its restriction
to \(B_{\mathrm{reg}}\), so a differentiated section extends
uniquely across the subset of codimension at least two.
The operator obeys
\[
             \mathcal D_z(us)
                    =\xi_z(u)s+u\mathcal D_z(s).
\]
Pull it to the chosen neighborhood in \(T_l\) using the lifted
symbol \(\widetilde v_0\): in a local pulled-back frame its rule is
\begin{equation}\label{iso:connection-operator}
 \widetilde{\mathcal D}_z(h\,p_l^*s)
       =\widetilde v_0(h)\,p_l^*s
                   +h\,p_l^*(\mathcal D_zs).
\end{equation}
The Leibniz identity and
\(\widetilde v_0(p_l^*u)=p_l^*(\xi_z u)\)
make this definition independent of the frame.

On an overlap, both operators
\(\widetilde{\mathcal D}_z\) have the same symbol.
They agree on the dense open where \(p_l\) is \'{e}tale and
both varieties are smooth: there the differential of \(p_l\)
identifies the pulled-back canonical tensors with canonical
tensors upstairs, and both operators are Lie differentiation
by the same field \(\widetilde v_0\).
Their difference is an \(\cO_{T_l}\)-linear endomorphism of the
line bundle \(\mathcal M\); its coefficient is holomorphic and
zero on a dense open, hence is zero.
Thus the local operators glue on a neighborhood of \(C\).
This argument compares the operators upstairs and does not
require a single descended vector field on a neighborhood of
all of \(S_0\).

Tangency of \(\widetilde v_0\) to the reduced curve \(C\) implies
that the glued operator preserves
\(\mathcal I_C\mathcal M\), by
\eqref{iso:connection-operator}.
It therefore induces a first-order operator
\[
             D:\mathcal M|_C\longrightarrow\mathcal M|_C
\]
with symbol the nowhere vanishing field
\(\widetilde v_0|_C\).
Let \(\tau\) be the holomorphic one-form on the elliptic curve
\(C\) satisfying \(\tau(\widetilde v_0)=1\).
Then
\[
                           \nabla s=D(s)\otimes\tau
\]
is a holomorphic connection on \(\mathcal M|_C\).

A line bundle with a holomorphic connection on a smooth projective
curve has degree zero, a classical connection obstruction
\cite{Atiyah1957}.  One elementary proof takes a nonzero
meromorphic section \(s\): the meromorphic one-form
\(\nabla s/s\) has at each point the residue
\(\operatorname{ord}_q(s)\), and the residue theorem gives
\(\deg\operatorname{div}(s)=0\).
Consequently
\[
                         \deg(\mathcal M|_C)=0.
\]
On the other hand, \(-K_B\) is ample as a \(\Q\)-Cartier divisor,
and \(p_l|_C\) is finite onto the positive-dimensional curve
\(S_0\).  Hence
\[
             \deg(\mathcal M|_C)
               =r\,\deg\bigl((p_l|_C)^*K_B\bigr)<0.
\]
This contradiction excludes the nonquotient germ and proves
Proposition~\ref{iso:exclusion}.
\end{proof}
